% Generated draft: requires mathematical and source review.
\documentclass[11pt]{article}
\usepackage[margin=25mm]{geometry}
\usepackage{fontspec}
\setmainfont{Latin Modern Roman}
\usepackage{amsmath,amssymb,mathtools,unicode-math}
\setmathfont{Latin Modern Math}
\usepackage{enumitem,needspace,xurl,hyperref}
\setlength{\parindent}{0pt}
\setlength{\parskip}{5pt}
\setlength{\emergencystretch}{3em}
\newcommand{\N}{\mathbb{N}}
\newcommand{\Z}{\mathbb{Z}}
\newcommand{\Q}{\mathbb{Q}}
\newcommand{\R}{\mathbb{R}}
\newcommand{\C}{\mathbb{C}}
\newcommand{\F}{\mathbb{F}}
\newcommand{\A}{\mathbb{A}}
\newcommand{\PP}{\mathbb{P}}
\newcommand{\E}{\mathbb{E}}
\newcommand{\eps}{\varepsilon}
\newcommand{\dd}{\,\mathrm{d}}
\begin{document}
\section*{UnsolvedMath ID: 30003253}
\textbf{Source problem code:} catalogue entry
\subsection{Definitions and mathematical statement}
Let \(X\) be a smooth projective surface over \(\C\), let \(\pi:X\to\PP^2_{\C}\) be a finite morphism with branch divisor \(B\), and fix \(g\geq0\). For the following formulation, a curve means a reduced irreducible curve. Its geometric genus is the genus of its smooth projective normalization. Set
\[
\mathcal C_g(X)=\{C\subset X:C\text{ is a curve of geometric genus at most }g\}.
\]
The displayed cohomology group in the source requires \(L\) to be an invertible sheaf on \(\PP^2_{\C}\). Under this interpretation, write \(L\simeq\mathcal O_{\PP^2}(d)\) and define its degree to be \(d\). Define
\[
\mathcal A_g=\left\{d\in\Z:\begin{array}{l}\text{there exists }0\ne s\in H^0(\PP^2,\mathcal O_{\PP^2}(d))\\\text{such that }s|_{\pi(C)_{\mathrm{red}}}=0\text{ for every }C\in\mathcal C_g(X)\end{array}\right\}.
\]
The minimum-degree hypothesis presupposes \(\mathcal A_g\ne\varnothing\). When this holds, let \(d_{\min}=\min\mathcal A_g\) and let \(L\simeq\mathcal O_{\PP^2}(d_{\min})\).

For a positive integer \(r\), \(S^r\Omega^1_{\PP^2/\C}\) denotes the \(r\)-th symmetric power of the cotangent sheaf. Given
\[
\omega\in H^0\!\left(\PP^2,L\otimes S^r\Omega^1_{\PP^2/\C}\right),
\]
and a reduced irreducible curve \(D\subset\PP^2\), let \(\nu_D:\widetilde D\to D\hookrightarrow\PP^2\) be its normalization followed by inclusion. Using the usual differential interpretation, call \(D\) \(\omega\)-integral when the image of \(\nu_D^*\omega\) under the natural map to
\[
H^0\!\left(\widetilde D,\nu_D^*L\otimes(\Omega^1_{\widetilde D/\C})^{\otimes r}\right)
\]
is zero.

The supplied clean statement asks whether the minimum-degree hypothesis guarantees a section \(\omega\) in the displayed cohomology group satisfying
\[
\begin{aligned}
& D\text{ is }\omega\text{-integral}
&&\text{for every irreducible component }D\text{ of }B,\\
& \pi(C)_{\mathrm{red}}\text{ is }\omega\text{-integral}
&&\text{for every }C\in\mathcal C_g(X).
\end{aligned}
\]
This remains a qualified formulation: the source does not specify the value, range, or quantifier for \(r\), and does not explicitly require \(\omega\ne0\). The positive-order definition above does not resolve those omissions. The original statement instead asks which hypotheses involving \(X\), \(B\), and \(L\) guarantee an appropriate differential and make a low-genus curve method work; it does not specify that method's required conclusion. The supplied background also mentions completeness of an integral-curve list and a higher-dimensional extension without formulating either target. These broader requests cannot be identified precisely with the clean existence question from the supplied record alone.
\subsection{Short English statement}
For a smooth projective complex surface admitting a finite map to \(\PP^2\), consider a plane line bundle of smallest degree having a nonzero section vanishing on every image of a curve of geometric genus at most \(g\). The supplied question asks for a twisted symmetric differential whose integral curves include these images and all branch components. The order of the differential, its required nonvanishing, and the broader method and extension targets remain unspecified.
\subsection{Sources}
\begin{itemize}
\item Supplied local UnsolvedMath registry record, permanent ID \(30003253\), exact source problem code \texttt{catalogue entry}, title ``Differential Criteria for Low-Genus Curves on Surfaces,'' in the Oberwolfach Reports Open Problems collection. The record gives proposed year 2016 and names no proposer.
\item Supplied report citation: \emph{Definability and Decidability Problems in Number Theory} (2017), Oberwolfach Reports, pp.~2793--2866, DOI \texttt{10.4171/owr/2016/49}. \url{https://doi.org/10.4171/owr/2016/49}. The supplied background describes these reports as records of mathematical workshops and open-problem discussions.
\item Supplied literature source: Garc\'ia-Fritz, \emph{Curves on surfaces}. \url{https://qspace.library.queensu.ca/bitstreams/007eb1de-ecf4-4747-a756-f7cac789ba56/download}.
\item Supplied literature source, labeled ``Oberwolfach Reports 2017'': \url{https://ems.press/content/serial-article-files/46652}.
\item Literature assessment dated exactly \texttt{2026-08-21}: the supplied curation labels the problem ``partially solved'' and states that surface-level criteria for completeness of the integral-curve list are known, while the general hypotheses and higher-dimensional extension remain unresolved. This is supplied dated curation, described in the record as surface-level literature triage, and has not been independently verified here. Its time-sensitive status and cited claims require independent checking.
\item The registry supplies the statement label \texttt{corrected\_verified} and reports that the clean statement was checked against the official Oberwolfach report and preserves its hypotheses, notation, and open target. This is a supplied verification claim, not an independent verification by this draft.
\end{itemize}
\end{document}
